\chapter{Introduction} \label{chapter:introduction}

\label{interface_importance} Material interfaces underpin the performance and reliability of advanced solid-state systems~\cite{Di_Liberto2022-gz}.
In semiconductor devices, they govern charge transport, carrier confinement, defect formation, and phase stability; in energy conversion and storage technologies, they regulate ion exchange kinetics, chemical compatibility, and thermal management~\cite{Maji2021-kc, Di_Liberto2022-gz, Cervantes2025-ze}.
In emerging heterostructures, interfacial interactions induce effects not found in the bulk, such as Schottky barriers or phase stabilisation~\cite{Low2025-on}.
As device dimensions shrink and surface-to-volume ratios increase, interfacial effects dominate performance, underscoring the importance of modelling interfaces with quantitative tools~\cite{Butler2019-ws, Leung2020-tv, Gerber2023-yd}.

\label{interface_foundation} Material interfaces delineate the boundary where two distinct solids meet, whether as heterostructure junctions, grain boundaries, or thin films~\nocite{}.
These exhibit unique atomic, electronic, and thermal characteristics that set them apart from the bulk phases~\cite{Gerber2023-yd, Di_Liberto2022-gz}.
Local features such as strain fields, atomic reconstruction, and elemental intermixing create modified bonding environments and electronic states, profoundly influencing the mechanical, transport, and optical behaviour of the composite~\cite{Khot2025-hh, Leitherer2023-ew, Maji2021-kc}.
In micro- and nanoscale systems, interfaces largely determine electronic and mechanical performance~\cite{Leung2020-tv}.

\newpage\label{material_interface_definition} In this research, material interfaces are studied predominantly through ab initio methods, particularly density functional theory (DFT) and new machine-learned interatomic potentials (MLPs)~\nocite{}.

Band offsets at interfaces, often deviating from predictions by simple models like Anderson's rule, are governed by local charge redistribution, dipole formation, and interface-specific states~\nocite{}.
The emergence of new electronic phases at interfaces, such as Ohmic contact formation or different band alignments, necessitates accurate modelling beyond bulk considerations~\nocite{}.

Several methodologies have been developed to address the challenges of constructing realistic interface models~\nocite{}.
The ARTEMIS tool enables high-throughput generation of interface structures through lattice matching and surface termination selection, allowing systematic exploration of the complex interfacial phase space~\cite{Taylor2020-nl}.
Similarly, the RAFFLE method uses statistical learning and void-filling algorithms to construct interface phases that are energetically favourable and representative of plausible morphologies~\cite{raffle_paper}.

Interfaces also significantly influence phonon transport~\cite{Khot2025-hh}.
They can reduce thermal conductivity compared to bulk values; an effect exploitable in thermoelectric applications~\cite{needed}.
Interfacial scattering, acoustic impedance mismatch, and atomic-scale disorder all contribute to this behaviour, reinforcing the interface's role as a functional component rather than a passive boundary~\cite{Hepplestone2010-vw}.
Thus, in the context of this research, a material interface is not merely a geometric boundary but an emergent, functionally distinct region~\nocite{}.
It governs critical physical properties, demands precise structural modelling, and enables advances in material innovation and device engineering~\nocite{}.

\newpage\label{interface_classification} Material interfaces can be broadly categorised by their crystallographic orientation, dimensionality, and degree of lattice coherence~\nocite{}.
These classifications are instrumental in understanding interfacial phenomena and in designing simulation protocols that are physically realistic and computationally tractable~\cite{Leung2020-tv, Di_Liberto2022-gz, Gerber2023-yd}.

\label{interface_coherence} Interfaces can also be distinguished by the degree of lattice match across the boundary:~\nocite{} % perhaps this would be a better place to cite Davies2024-ju?

\begin{itemize}
    \item Coherent interfaces exhibit a perfect registry between atomic planes, typically occurring when lattice mismatch is negligible~\nocite{cited_below}.
    These interfaces preserve the translational symmetry across the junction and are often found in homoepitaxial or pseudomorphic systems~\cite{Davies2024-ju}. % citation might be better in another location
    \item Semi-coherent interfaces accommodate lattice mismatch through periodic misfit dislocations~\nocite{cited_below}.
    These dislocations relieve interfacial strain while maintaining partial registry, allowing larger domains of coherency interspersed with localised defects~\cite{Davies2024-ju}. % citation might be better in another location
    \item Incoherent interfaces display no periodicity match across the interface, leading to structurally disordered junctions~\nocite{cited_below}.
    These are often found in polycrystalline or amorphous/crystalline composites and can act as strong phonon or electron scattering centres~\cite{Davies2024-ju}. % citation might be better in another location
\end{itemize}

\label{chemical_contrast_hetero_homo} Chemical composition also plays a contrasting and defining role~\nocite{}.

\begin{itemize}
    \item Heterojunctions involve distinct materials with differing electronegativity, band gaps, and ionic/covalent character~\nocite{}.
    Heterojunctions are central to semiconductor devices, solar cells, and catalysis, often giving rise to emergent interfacial dipoles, states, or reconstructed bonding motifs~\nocite{}.

    A heterostructure is a structure composed of multiple heterojunctions~\nocite{}.
    \item Homojunctions are interfaces within a single material where the crystal lattice remains continuous, but local properties differ, such as variations in dopant concentration or applied strain, resulting in regions with distinct electronic or mechanical behaviour~\nocite{}.

    As with heterojunction vs structure, homostructure is a structure composed of multiple homojunctions~\nocite{}.
\end{itemize}

\label{interface_effects_2D_oxides} % Material interfaces influence a wide array of physical and electronic properties, often determining the functional behaviour of composite systems and devices~\nocite{}.

These effects are particularly pronounced in systems involving 2D materials, complex oxides, or strongly correlated electron systems~\cite{Leitherer2023-ew}.

\label{band_alignment_interfaces} % Perhaps the most critical electronic property influenced by interfaces is the band alignment between materials~\nocite{}.
% Misalignment of conduction and valence band edges creates energy barriers or wells that determine whether an interface behaves as a Schottky barrier, Ohmic contact, or heterojunction~\nocite{}.
% These alignments are crucial for device operation in field-effect transistors, diodes, and photovoltaics~\cite{Chagarov2015-ol, Schusteritsch2015-lw, Low2025-on}. % bring back?
% Interfacial dipoles, charge transfer, and electronic reconstruction can shift the effective band edges, often invalidating simple rules like Anderson's model~\cite{Chagarov2015-ol, Gerber2023-yd, Low2025-on}. % bring back?

\label{interface_states} Interfaces often host localised states arising from atomic mismatch, bonding reconstruction, or impurity segregation~\nocite{}.
These can act as traps or recombination centres, significantly affecting carrier's lifetimes and mobility~\cite{Hepplestone2010-vw, Maji2021-kc, Davies2024-ju}.

\label{dielectric_interface_permittivity} The dielectric properties of materials can be strongly altered at interfaces~\nocite{}.
In systems like $\ce{CaCu3Ti4O12}$, colossal permittivity arises not from the bulk crystal but from insulating grain boundaries and a semiconducting interior, exemplifying the Maxwell-Wagner effect~\cite{needed}.
Interface-induced polarisation, charge accumulation, and ionic displacements all contribute to enhanced dielectric responses~\cite{Leung2020-tv, Low2025-on}.

\label{thermal_transport_interfaces} Thermal transport across interfaces is typically suppressed due to phonon scattering, mass disorder, and acoustic impedance mismatch~\cite{needed}.
This results in thermal boundary resistance (Kapitza resistance), which strongly influences heat flow in thermoelectric and microelectronic devices~\cite{Hepplestone2010-vw, Leung2020-tv, Khot2025-hh}.
In heterostructures, the interface can significantly scatter phonons, reducing thermal conductivity below that of either constituent~\cite{Di_Liberto2022-gz, Davies2024-ju}.

\label{interface_stabilise_metastable_phases} Interfaces can stabilise new or metastable phases not accessible in bulk~\cite{needed}. % hepplestone 2010 i think?
For instance, the interface between $\ce{BaTiO3}$ and $\ce{SiO2}$ has been shown to promote the spontaneous formation of the fresnoite phase $\ce{Ba2TiSi2O8}$, altering both electronic and mechanical response~\cite{needed}. % hepplestone and ned's student paper
A more nuanced example is the $graphene|\ce{MgO}$ interface: while monolayer MgO confined between graphene layers often adopts a rocksalt-like structure, this is not simply a case of one phase being energetically favoured~\nocite{}. % hepplestone and ned's student paper
Instead, the local environment, including charge transfer and encapsulation, stabilises the rocksalt phase relative to the expected hexagonal phase in unsupported monolayers~\nocite{}. % hepplestone and ned's student paper
However, experimental and theoretical studies suggest that a mixed or hybrid phase may also emerge with phase stability dependent on thickness and interface strain~\cite{needed}. % hepplestone and ned's student paper

\newpage\label{interface_mechanical_properties} Mechanical properties such as bulk modulus, shear strength, and adhesion energy are modified at interfaces due to altered bonding environments and lattice mismatch~\nocite{}.
These variations are important in the design of composite materials, flexible electronics, and nanoindentation studies~\cite{Zhang2016-lt, Leung2020-tv, Davies2024-ju}.

\label{interface_optical_effects} Interfaces affect the optical absorption spectrum by introducing new electronic states and modifying the joint density of states~\nocite{}.
Strain-induced effects at the interface can shift bandgap and optical properties, which are useful for photodetectors and photovoltaic devices~\cite{Chagarov2015-ol, Low2025-on}.

\label{interface_summary} Accurate control and prediction of material interfaces are vital for the design of next-generation electronic, optoelectronic, and thermal devices~\cite{Low2025-on}.

\label{interface_functional_regions} % Interfaces are not merely passive boundaries between materials; they have an active impact on their surrounding environment, functional regions that frequently control and even dominate the physical properties of the system.
% As materials and devices are engineered at ever smaller scales, the role of interfaces becomes more pronounced~\cite{Chagarov2015-ol}.
% In conventional bulk systems, interfaces may be limited to grain boundaries or interphase regions~\cite{Leitherer2023-ew}.
% In nanomaterials, layered heterostructures, and device-scale systems, interfacial regions can constitute a non-negligible fraction of the total volume or surface area~\cite{Leitherer2023-ew}.
% Their influence is therefore not ancillary but foundational to system performance~\nocite{}.

\label{interface_stability_challenge} Yet despite their importance, predicting which interfaces are energetically favourable remains a significant challenge~\nocite{}.
% Interface stability is not solely determined by bulk lattice matching, but also by local bonding, chemical inhomogeneity, and atomic-scale reconstruction~\nocite{}.
The structural space is vast: each pair of surfaces can be stacked with arbitrary lateral shifts, terminations, and alignments, resulting in a combinatorial explosion of configurations~\cite{Zhang2016-lt, Di_Liberto2022-gz, Gerber2023-yd}.
Strain relaxation, interlayer interactions, and charge transfer further complicate the energy landscape~\nocite{}.
Accurately assessing favourability across this space demands a balance between resolution and scale that conventional methods struggle to achieve~\cite{Gerber2023-yd, Deng2023-tg, Low2025-on}.

\label{dft_scalability_bottleneck} First-principles techniques such as DFT remain the standard for computing accurate interfacial energies and relaxations~\cite{Cervantes2025-ze}.
However, their high computational cost, scaling cubically $O(n^3)$ with system size, renders exhaustive sampling intractable for realistic sizes structures, which often comprise hundreds to thousands of atoms~\cite{Schutt2019-sz, Butler2019-ws, Rosen2022-pw}.
This bottleneck precludes broad screening efforts, especially when exploring multiple chemistries, orientations, and stacking arrangements~\nocite{}.
As such, DFT might be better suited to final validation and a confirmation of datasets evaluated in a less computationally expensive manner~\nocite{}.

\newpage\label{:mlp_screening} The present project explores the use of MLPs for scalable interface screening~\nocite{}.
MLPs are data-driven models trained to reproduce DFT-level energies and forces, but at orders of magnitude lower cost~\cite{Wang2024-tn}.
In particular, this work focuses on the MACE framework; an equivariant message-passing architecture that preserves rotational symmetries and captures local atomic interactions with high fidelity~\nocite{}.
By using MACE to rapidly evaluate large numbers of candidate interfaces, it becomes possible to rank structures by predicted stability and prioritise low energetic configurations for subsequent DFT validation~\cite{Batatia2022-xz}.

\label{hierarchical_modelling_workflow} This approach supports a hierarchical modelling workflow in which DFT and MLPs are used in tandem: MLPs for breadth, and DFT for depth~\nocite{}.
It enables rigorous testing of whether MLPs can replicate DFT-derived favourability rankings across a range of systems, including phase boundaries (e.g. $\ce{Si}|\ce{Si}$) and heterostructures (e.g. $\ce{Si}|\ce{Ge}$, $\ce{Si}|\ce{C}$)~\nocite{}.
% For the context of this study, I ignore the semicontuctor requirement in the definition of heterostructure~\nocite{}.
Moreover, it lays the foundation for further improvements through techniques such as delta-learning, structure-based prediction models, and surface-to-interface generalisation~\nocite{}.

\label{report_structure} The remainder of this report is structured as follows~\nocite{}:

\begin{itemize}
    \item \textbf{Chapter~\ref{chapter:background}} reviews the theoretical and computational background for interface modelling~\nocite{}.
    It defines material interfaces, explores interfacial energetics and physical effects, and surveys applications across electronic, catalytic, and energy systems~\nocite{}.
    The chapter also introduces DFT, MLPs, and recent algorithmic tools for generating and evaluating heterostructures~\nocite{}.
    \item \textbf{Chapter~\ref{chapter:computational_investigation}} evaluates the use of MACE as a surrogate model for efficiently screening interface structures prior to DFT~\nocite{}.
    It benchmarks MACE vs DFT energetics across over roughly 680 heterostructures, analyses rank correlations, and identifies system-specific performance trends that inform future dataset construction for predictive modelling~\nocite{}.
    \newpage\item \textbf{Chapter~\ref{chapter:next_steps}} outlines the design and evaluation of a Machine-Learned Interface (MLI) model trained to predict interface favourability~\nocite{}.
    It compares three training strategies baseline MACE, interface-specific fine-tuning, and delta-learning correction, and introduces ranking metrics to quantify predictive performance~\nocite{}.
    The chapter also explores how structural asymmetries, surface orientation, and chemical diversity impact model generalisation~\nocite{}.
    \item \textbf{Chapter~\ref{chapter:conclusions}} reflects on key findings across materials systems, highlighting the strengths and limitations of current MLPs in replicating DFT interface rankings~\nocite{}.
%    It emphasises the importance of rank-based metrics over absolute errors and advocates for hybrid workflows combining fast MLP screening with selective DFT correction to guide future materials discovery~\nocite{}.
\end{itemize}

\label{framework_objective} Through this integrated modelling framework, the project seeks to reduce the cost and uncertainty of interface prediction while preserving the fidelity needed for scientific insight and technological relevance~\nocite{}.

\newpage\section*{Motivation for the Research}\label{section:motivation}

\label{nonideal_interface_formation} Interfaces rarely form under ideal conditions~\nocite{}.
Instead, they emerge through complex, often non-equilibrium processes such as high-temperature annealing, epitaxial growth, and chemically driven surface reconstructions~\nocite{}.
These dynamics frequently trap systems in metastable configurations or generate local disorder, resulting in atomic arrangements that deviate significantly from those predicted by ideal lattice-matching rules~\nocite{}.
Such complexity challenges efforts to model interface behaviour using conventional approaches alone~\cite{Hepplestone2010-vw, Leung2020-tv, Di_Liberto2022-gz, Gerber2023-yd}.

\label{dft_scaling_limitations} While DFT provides a reliable route to total energy prediction and structural relaxation, its unfavourable scaling with system size ($\mathcal{O}(N^3)$ in electron count) severely limits its use in high-throughput workflows~\nocite{}.
In particular, DFT becomes prohibitive when applied to interface structures exceeding a few hundred atoms, as required to capture long-range strain, multi-layer reconstruction, or compositional variation~\cite{Butler2019-ws, Gerber2023-yd}.
Moreover, standard DFT approximations, such as LDA exchange-correlation functionals, often struggle to predict key interfacial phenomena in layered systems~\nocite{}.
For example, accurate band-alignment at a heterojunction determines charge-injection barriers and device turn-on voltages; and weak van der Waals interactions between two-dimensional layers govern adhesion, interlayer spacing, and electronic reconstructions~\nocite{}.
Failure to capture these effects can lead to qualitatively incorrect predictions of interface stability and function and requires a higher level of DFT~\cite{Gerber2023-yd, Low2025-on}.

\label{mlp_overview} MLPs offer a promising route to overcoming these limitations~\nocite{}.
By learning to emulate DFT-level energies and forces, MLPs like MACE enable large-scale structure relaxations and energy evaluations at dramatically reduced computational cost~\cite{Khot2025-hh, Wang2024-tn}.
This opens up the exploration of far broader configurational spaces, including low-symmetry interfaces and extended defects, that would otherwise be intractable using DFT alone~\nocite{}.
However, this benefit is only realised if the models generalise well across diverse interfacial environments~\nocite{}.
Achieving DFT-comparable accuracy in such settings requires a representative, high-fidelity training set: any gaps or systematic biases in the reference data will be reflected in the potential's predictions~\cite{Schutt2017-ge, Batzner2022-fd, Deng2023-tg}.

\label{motivation_mlp_dft} % The primary aim of this report is to evaluate whether MLPs can reproduce the DFT-derived ranking of interface favourability, here quantified by formation energy, across a set of group 14 heterostructures~\nocite{}.
%While the overarching PhD objective is to develop a Machine-Learned Interface (MLI) model that predicts interface energies directly from the constituent slabs (bypassing explicit interface construction), this study focuses on rapid MLP-assisted screening to identify the most promising candidates for eventual full DFT ISIF3 relaxations~\nocite{}.
Chapter~\ref{chapter:computational_investigation} shows that MLPs often preserve relative DFT energy rankings, particularly within chemically similar systems, with a few edge cases like pure-\ce{Si} interfaces (grain boundaries)~\nocite{}.
These discrepancies underscore the need for targeted refinement, motivating the delta-learning approaches presented in Chapter~\ref{chapter:next_steps}, which correct systematic MLP-DFT residuals via supervised learning on interface-specific error distributions~\nocite{}.

\label{project_motivation} Ultimately, this project is motivated by the need for more efficient and predictive tools to accelerate the design of functional heterostructures~\nocite{}.
By embedding MLPs within automated screening workflows and augmenting them with delta-learning corrections and structure-based heuristics, it becomes possible to bridge the gap between physical realism and computational tractability~\nocite{}.
In doing so, this research contributes to the broader aim of enabling scalable, data-driven exploration of complex interfacial systems; particularly in regimes where bulk-derived heuristics break down~\nocite{}.

\newpage\section*{Research Objectives} \label{section:research_objectives}

\label{objectives_mli_goal} This PhD aims to develop a Machine-Learned Interface (MLI) model that, given two input materials; whether slabs, bulks, or combinations, can generate realistic and physically plausible interface structures between them~\nocite{}.
The purpose of this model is not to directly predict interfacial energies or properties, but to output candidate geometries that are likely to be realisable and energetically reasonable, enabling downstream evaluation via existing MLPs or DFT~\nocite{}.

\label{objectives_mli_pipeline} \textbf{The ultimate goal is to accelerate interface discovery by replacing brute-force combinatorial stacking with an intelligent generative model~\nocite{}.}
By learning patterns in crystallography, registry, and bonding across known interfaces, the MLI model will propose viable configurations that reduce the need for exhaustive enumeration and manual filtering~\nocite{}.

\label{objectives_mli_extended_use} Such a tool could eventually support applications beyond interface prediction, including surface termination prediction under reactive gas atmospheres, simulation of crystal growth through layer-by-layer stacking, and real-time assessment of growth defects during slab assembly~\nocite{}.

\label{objectives_dataset_building} To support development of the MLI model, the first phase of this research focuses on constructing a large, diverse dataset of heterostructure interfaces~\nocite{}.
Since direct generation and relaxation of thousands of interfaces using DFT would be prohibitively expensive, machine-learned interatomic potentials (MLPs), specifically the MACE framework, are evaluated for their ability to replicate DFT-derived rankings of interface favourability~\cite{Batatia2022-xz}.

\label{objectives_screening_workflow} % ---------- %

If MLPs can reliably identify promising structures, even without full relaxation, they can serve as a low-cost screening tool to filter candidate interfaces for inclusion in the training set~\nocite{}.

This enables the creation of a scalable data pipeline where realistic interfaces are generated via tools like ARTEMIS~\cite{Taylor2020-nl},
screened using MLPs, and selectively validated with DFT~\nocite{}.

\label{objectives_dft_validation} Once promising interfaces are identified, a subset will be fully relaxed using DFT to provide high-fidelity data for MLI training~\nocite{}.
These DFT-relaxed structures will serve as the reference for extracting physically meaningful descriptors, such as surface orientation, lattice mismatch, interfacial strain, atomic registry, and coordination environments, that can be used as features for the MLI model~\cite{Gouveia2025-fpml}.

\label{objectives_dft_mlp_comparison} Although initial results show strong rank correlation between MLP and DFT energies (e.g., $\rho \approx 0.98$ for Sn|Ge), this agreement may degrade in more complex systems or larger interface models.
Full DFT validation remains essential to calibrate and benchmark the workflow, particularly for low-symmetry or poorly coordinated interfaces where MLP generalisation may falter.

\label{objectives_relaxation_scaling} To further extend this validation, rankings based on MLP-relaxed geometries will be compared to those from DFT-relaxed structures~\nocite{}.
If similar rank orderings are recovered across both relaxation methods, it would support the use of MLPs for handling much larger and more realistic interfacial systems, ones that are computationally infeasible to relax using DFT alone~\nocite{}.
This opens the possibility of generating structurally diverse, high-volume datasets beyond the reach of traditional ab initio approaches~\nocite{}.

\label{objectives_inverse_design} % ---------- %

A longer-term extension of this research, contingent on the success of the main MLI model and access to further computational resources, is the development of an inverse interface design framework~\nocite{}.
Rather than proposing interfaces from input slabs alone, this model would take target properties or desired interfacial features, such as low strain, specific registry types, or stability thresholds, and suggest plausible $A|B$ configurations that meet those criteria~\nocite{}.

\label{objectives_inverse_limitations} Such a tool could support experimental design by guiding researchers toward synthesizable interfaces that meet functional requirements~\nocite{}.
However, training such a model would require a much broader dataset, including properties computed from full DFT relaxations on MLI-generated structures, something only feasible with additional funding and computational time~\nocite{}.

\label{objectives_mlp_refinement} To improve the accuracy of interface energy predictions, particularly for systems where MACE performs poorly (e.g. strained silicon-rich structures or cross-material environments), this project will explore refinement strategies for the MLP itself~\nocite{}.
One approach is delta-learning: training a secondary model to correct systematic energy differences between MACE and DFT, using DFT-relaxed structures as a reference ~\cite{Pitfield2025-ah}.

\label{objectives_mlp_tuning} Another strategy is direct fine-tuning or retraining of MACE on an expanded dataset that includes interface-specific configurations, such as strained slabs, low-symmetry orientations, or cross-material bonding~\nocite{}.
This is motivated by known failure modes in MLIPs, including poor generalisation to novel bonding environments, and underprediction of surface and interface energies~\cite{Leung2020-tv, Uhrin2025-ba}.
These efforts will aim to mitigate known risks; such as extrapolation errors, by enriching the training distribution with interfacial scenarios that challenge model stability~\nocite{}.
By explicitly introducing diverse bonding motifs and strain conditions, the model's predictive performance can be made more robust across real-world interface configurations, reducing the likelihood of catastrophic errors in high-throughput screening~\nocite{}.

\label{objectives_failure_modes} A further objective of this research is to identify and mitigate failure modes in machine-learned potentials when applied to interfacial systems~\nocite{}.

Despite strong average performance, modern MLIPs such as MACE exhibit known limitations, including potential energy surface (PES) softening, poor extrapolation to strained or defective configurations, and underprediction of surface energies~\cite{Batzner2022-fd}.

\label{objectives_failure_diagnostics} This work will systematically assess when and why these failures occur, starting with Group IV interfaces as test cases.
Specific attention will be paid to bonding asymmetries, coordination anomalies, and cases where poor generalisation leads to incorrect rank ordering or unrealistic relaxations.
Where needed, refinement strategies such as delta-learning, fine-tuning, and active learning will be deployed to correct these behaviours~\cite{needed}.
The goal is not only to improve model accuracy, but also to build a deeper understanding of where and why MLPs break down when applied to out-of-distribution geometries; informing the development of more robust, uncertainty-aware workflows for interface prediction~\nocite{}.

\newpage\label{objectives_structure_insight} In parallel with predictive modelling, the project seeks to uncover the structural and chemical features that determine interfacial stability~\nocite{}.
One key line of enquiry is whether surface favourability correlates with interface favourability; for example, whether low-energy surfaces tend to form stable interfaces, or whether compensating effects between unfavourable surfaces can yield unexpectedly favourable interfaces~\nocite{}.

\label{objectives_structure_analysis} Further investigations will explore how stability varies across crystal types (e.g. fcc, hexagonal, tetragonal), material classes (e.g. semiconductors, metals, insulators), and bonding mechanisms (e.g. covalent, ionic, metallic)~\nocite{}.
The influence of valence orbitals; such as $s$ versus $p$ dominance, will also be considered, particularly in how they affect registry, relaxation, and bonding across the interface~\nocite{}.

\label{objectives_summary} Together, these objectives form a cohesive research trajectory: from benchmarking machine-learned potentials and constructing curated datasets, to training a generative interface model (MLI) and exploring its extension toward inverse design~\nocite{}.
The project aims not only to reduce the computational cost of interface exploration, but also to uncover structural principles and predictive strategies that generalise across materials systems~\nocite{}.

\label{objectives_impact} By integrating machine learning with crystallographic modelling and ab initio validation, this work supports the development of practical tools for scalable, data-driven interface discovery in electronic, optoelectronic, and catalytic materials~\nocite{}.

% ---------- %

\newpage\section*{Scientific and Practical Significance} \label{section:scientific_and_practical_significance}

\label{interfaces_nanoscale_behaviour} Material interfaces are central to the behaviour of functional solids, especially at the nanoscale where surface and interfacial effects dominate bulk properties~\cite{Chagarov2015-ol, Leitherer2023-ew, Khot2025-hh, Cervantes2025-ze}.
In semiconductors, they define charge transport pathways by setting potential barriers and band offsets~\cite{Low2025-on}.
In batteries, they regulate ionic conductivity, chemical stability, and degradation kinetics~\cite{Leung2020-tv}.
In photovoltaics and quantum devices, interfacial dipoles, tunnelling channels, and confinement effects underpin key device operations~\cite{Low2025-on, Cervantes2025-ze}.

\label{interfacial_electronic_shifts} Subtle differences in stacking, coherence, or atomic reconstruction can shift critical electronic properties by several electronvolts; altering conductivity, carrier lifetimes, or activation barriers~\nocite{}.
These shifts emerge from a complex interplay of interfacial strain, chemical bonding, and electronic redistribution, and are sensitive to both global structure and local environment~\cite{Batzner2022-fd, Gerber2023-yd, Davies2024-ju}.
As such, predicting interface stability and functionality from atomic configuration alone remains a demanding task~\nocite{}.

\label{dft_scaling_cost} While DFT offers a robust ab initio framework for computing interfacial energies and electronic properties, its computational cost scales poorly with system size~\cite{Schutt2019-sz, Leung2020-tv, Cervantes2025-ze}.
Realistic supercells, often comprising hundreds of atoms, are therefore difficult to study in large numbers~\nocite{}.
DFT is well-suited to final validation, but ill-suited to exhaustive exploration across configurational space~\nocite{}.

\label{mlp_workflow_integration} Machine learning (ML) models, particularly MLPs, provide a way forward~\nocite{}.
When trained on DFT data, MLPs can reproduce formation energies and atomic forces at orders-of-magnitude lower cost, making them suitable for broad screening applications~\nocite{}.
Their integration into automated workflows, using tools such as ARTEMIS and RAFFLE, enables systematic generation and evaluation of interface candidates across crystallographic orientations, stacking shifts, and chemical combinations~\cite{artemis_paper, raffle_paper}.

\newpage\label{ml_interface_dynamics} Importantly, the utility of ML does not stop at static energetics~\nocite{}.
Many applications require insight into how interfaces behave under external stimuli; electric fields, temperature gradients, or chemical exposure~\cite{Chagarov2015-ol, Butler2019-ws, Leung2020-tv, Khot2025-hh}.
More adaptive and scalable modelling approaches are required for emerging interfacial phases, such as field-tunable dielectric layers or metastable heterostructures~\nocite{}.

\label{project_contributions} This project contributes to that broader effort.
It assesses the ability of MLPs, specifically MACE, to recover DFT-derived favourability rankings, identifies where discrepancies arise, and proposes delta-learning corrections to address them~\nocite{}.
It also lays groundwork for predictive models that estimate interface properties from top and bottom slabs, enabling pre-screening before full interface construction~\nocite{}.
By focusing on group-IV semiconductor systems and both homostructure and heterostructure cases, this work targets a class of technologically relevant interfaces where modelling uncertainty remains high and conventional rules often fail~\nocite{}.

\label{advancing_interface_discovery} Ultimately, this research supports the advancement of data-driven, scalable methods for interface discovery; bridging the gap between computational feasibility and the physical fidelity required for next-generation materials design~\nocite{}.
